\documentclass[aspectratio=169,11pt]{beamer}
\usepackage[T1]{fontenc}
\usepackage[utf8]{inputenc}
\usepackage[brazil]{babel}
\usepackage{lmodern}
\usepackage{booktabs}
\usepackage{array}
\usepackage{amsmath}
\definecolor{Azul}{HTML}{17324D}
\definecolor{Verde}{HTML}{087F8C}
\definecolor{Cinza}{HTML}{566573}
\definecolor{Fundo}{HTML}{F0F5F8}
\definecolor{Terracota}{HTML}{B04B35}
\setbeamercolor{normal text}{fg=Azul,bg=white}
\setbeamercolor{structure}{fg=Verde}
\setbeamercolor{frametitle}{fg=Azul,bg=white}
\setbeamercolor{block title}{fg=Azul,bg=Fundo}
\setbeamercolor{block body}{fg=Azul,bg=Fundo}
\setbeamercolor{alerted text}{fg=Verde}
\setbeamerfont{frametitle}{size=\Large,series=\bfseries}
\setbeamertemplate{navigation symbols}{}
\setbeamertemplate{itemize items}[circle]
\setbeamersize{text margin left=8mm,text margin right=8mm}
\setlength{\parskip}{0pt}
\setbeamertemplate{itemize/enumerate body begin}{\setlength{\itemsep}{2mm}}
\newcommand{\secao}{LEITURA DE ARTIGO}
\newcommand{\fonte}[1]{\vfill\par{\fontsize{7}{8}\selectfont\color{Cinza} Fonte: #1\par}}
\setbeamertemplate{headline}{}
\setbeamertemplate{frametitle}{\vspace{1.5mm}{\fontsize{7}{8}\selectfont\color{Verde}\secao\par}\vspace{1.5mm}{\usebeamerfont{frametitle}\insertframetitle\par}\vspace{1.5mm}}
\setbeamertemplate{footline}{\leavevmode\hbox{\begin{beamercolorbox}[wd=.85\paperwidth,ht=3ex,dp=1.5ex,leftskip=8mm]{normal text}\fontsize{7}{8}\selectfont\color{Cinza} Welyab Paula\quad |\quad Anggoro et al. (2025)\end{beamercolorbox}\begin{beamercolorbox}[wd=.15\paperwidth,ht=3ex,dp=1.5ex,center]{normal text}\fontsize{7}{8}\selectfont\color{Cinza}\ifnum\value{framenumber}>14 Complemento\else\insertframenumber/14\fi\end{beamercolorbox}}}
\setbeameroption{hide notes}
\hypersetup{pdftitle={Ajuste contrastivo para classificação fiscal - Anggoro et al. (2025)},pdfauthor={Welyab Paula},pdfsubject={SBERT, MNR e classificação de códigos HS; apresentação de aproximadamente 13 minutos}}
\DeclareFontShape{T1}{lmss}{bx}{sl}{<->ssub*lmss/bx/n}{}
\DeclareFontShape{T1}{lmss}{bx}{it}{<->ssub*lmss/bx/n}{}
\DeclareMathSizes{10}{10.95}{8}{6}
\begin{document}
\renewcommand{\secao}{LEITURA DE ARTIGO}
% Slide 1; 30 segundos
\begin{frame}[t]
\normalfont\sffamily
{\fontsize{7}{8}\selectfont\color{Verde} LEITURA DE ARTIGO\par}
\par\vspace{3mm}
{\LARGE\bfseries Ajuste contrastivo para\\ classificação fiscal\par}
\par\vspace{1.5mm}
{\small\itshape Harmonized System Code Classification using Supervised Contrastive Learning with Sentence BERT and Multiple Negative Ranking Loss\par}
\par\vspace{1.5mm}
Angga Wahyu Anggoro, Padraig Corcoran,\\ Dennis De Widt e Yuhua Li\par
\par\vspace{1.5mm}
{\small Data Technologies and Applications, 59(2), 279--301, 2025\par}
\vfill
{\small Welyab Paula\quad |\quad PECS / UEMA\par}
\fonte{Anggoro et al. (2025), folha ORCA e resumo.}
\end{frame}
\note{\small Apresentar o estudo e sua pergunta: ajustar a representação das descrições comerciais melhora a previsão de códigos fiscais? O foco será o artigo de Anggoro; o artigo de Welyab fornece a motivação e a conexão com NF-e.\par [Sources] Anggoro et al. (2025), folha ORCA e resumo.}

\renewcommand{\secao}{MOTIVAÇÃO}
% Slide 2; 55 segundos
\begin{frame}[t]{O problema também aparece nas NF-e}
\normalfont\sffamily
\textbf{Descrições livres tornam a classificação fiscal difícil.}
\par\vspace{1.5mm}
\begin{columns}[T,onlytextwidth]
\begin{column}{.47\textwidth}
\textbf{Welyab et al.}\par\par\vspace{1.5mm}
\begin{itemize}
\item NF-e da SEFAZ-MA, em português.
\item Embeddings prontos e classificadores supervisionados.
\item Previsão de capítulos da NCM.
\end{itemize}
\end{column}
\begin{column}{.47\textwidth}
\textbf{Anggoro et al.}\par\par\vspace{1.5mm}
\begin{itemize}
\item Declarações de importação, em inglês.
\item Ajuste contrastivo do gerador de embeddings.
\item Previsão de códigos dentro de dois capítulos.
\end{itemize}
\end{column}
\end{columns}\par
\par\vspace{1.5mm}
\begin{block}{Pergunta central}
Aprender a representação no domínio melhora a classificação?
\end{block}
\fonte{Welyab et al., seções 3--4; Anggoro et al., resumo e seção 3.}
\end{frame}
\note{\small As duas pesquisas usam descrições reais, sujeitas a abreviações e ambiguidade. Em Welyab, os embeddings são gerados uma vez e os classificadores são treinados. Anggoro atualiza também o codificador. O ponto comum é texto, vetor e previsão de classe. O uso de rótulos declarados não comprova que o enquadramento fiscal está correto.\par [Sources] Welyab et al., seções 3--4; Anggoro et al., resumo e seção 3.}

\renewcommand{\secao}{DADOS}
% Slide 3; 45 segundos
\begin{frame}[t]{Duas bases reais de importação}
\normalfont\sffamily
\begin{center}
\renewcommand{\arraystretch}{1.35}
\begin{tabular}{@{}lrr@{}}
\toprule
& \textbf{Índia} & \textbf{Estados Unidos}\\
\midrule
Ano das transações & 2016 & 2018\\
Fonte & Zauba & Enigma\\
Transações no recorte & 66.522 & 58.003\\
Classes & \textbf{172} & \textbf{112}\\
\bottomrule
\end{tabular}\par
\end{center}
\par\vspace{1.5mm}
\begin{itemize}
\item Descrições em inglês; dados de declarações aduaneiras.
\item Recorte: capítulos \textbf{84 e 85} do Sistema Harmonizado.
\item Classes desbalanceadas: razão média relatada de \textbf{30:1}.
\end{itemize}
\fonte{Anggoro et al., seção 3.1 e Tabela I.}
\end{frame}
\note{\small São duas bases públicas, escolhidas também pelo alinhamento do idioma com BERT e DistilBERT em inglês. Os números da tabela são do recorte experimental, não do universo inteiro de cada fonte. Os capítulos abrangem máquinas e equipamentos mecânicos e elétricos. O artigo relata desequilíbrio entre classes majoritárias e minoritárias.\par [Sources] Anggoro et al., seção 3.1 e Tabela I.}

\renewcommand{\secao}{DADOS}
% Slide 4; 55 segundos
\begin{frame}[t]{Descrição é entrada; código é o rótulo}
\normalfont\sffamily
\textbf{Entrada:} descrição comercial do produto.\par
{\small Pode conter marca, modelo, capacidade, dimensões e número de série.}
\par\vspace{1.5mm}
\renewcommand{\arraystretch}{1.3}
\begin{tabular}{@{}lp{.70\textwidth}@{}}
\toprule
\textbf{Código} & \textbf{Descrição original (trecho)}\\
\midrule
84713010 & X2 LAPTOP LENOVO MODEL THINKPAD T430\\
85285100 & LED MONITOR, S22E360H, 21.5, INDIA \ldots\\
85423100 & INTEGRATED CIRCUIT LOGIC MONO ASTBL \ldots\\
\bottomrule
\end{tabular}\par
\par\vspace{1.5mm}
\textbf{Features dos classificadores:} 768 dimensões do embedding.\par
\par\vspace{1.5mm}
\alert{84 e 85 delimitam o recorte; não são as duas classes previstas.}
\fonte{Anggoro et al., seções 3.1 e 4.3.1; Tabela II.}
\end{frame}
\note{\small Cada linha é uma transação. O código serve de supervisão e não de entrada na inferência. Marca, peso e dimensões permanecem no texto, não são colunas de features separadas. A tarefa distingue 172 ou 112 códigos. O HS é internacional até seis dígitos; a base pode acrescentar extensões locais. A Tabela II exibe códigos de oito dígitos, que não devem ser chamados de NCM brasileira. O artigo não enumera aqui todas as classes nem documenta explicitamente uma padronização única do comprimento nas duas bases.\par [Sources] Anggoro et al., seções 3.1 e 4.3.1; Tabela II.}

\renewcommand{\secao}{CONCEITOS}
% Slide 5; 45 segundos
\begin{frame}[t]{SBERT transforma uma descrição em vetor}
\normalfont\sffamily
\textbf{Embedding:} representação numérica de um texto.\par
\par\vspace{1.5mm}
\begin{center}
Descrição do produto\\[1mm]
$\downarrow$\\[1mm]
\textbf{BERT ou DistilBERT}\quad +\quad \textbf{pooling}\\[1mm]
$\downarrow$\\[1mm]
Vetor com \textbf{768 dimensões}
\end{center}
\par\vspace{1.5mm}
\begin{itemize}
\item \textbf{SBERT:} arquitetura para representar textos inteiros.
\item \textbf{Pooling:} agrega os vetores dos tokens (média, máximo ou [CLS]).
\item O ajuste atualiza os parâmetros do codificador.
\end{itemize}
\fonte{Anggoro et al., seções 3.2 e 4.3.1.}
\end{frame}
\note{\small BERT gera representações dos tokens; pooling reúne essas representações em um vetor fixo para a descrição inteira. SBERT é um framework com codificador e pooling, não apenas o nome de um único modelo pronto. O artigo monta essa arquitetura com BERT ou DistilBERT pré-treinados; não treina um transformer do zero. DistilBERT é uma versão menor, usada para estudar o compromisso entre desempenho e tamanho.\par [Sources] Anggoro et al., seções 3.2 e 4.3.1.}

\renewcommand{\secao}{CONCEITOS}
% Slide 6; 80 segundos
\begin{frame}[t]{O mesmo código define pares positivos}
\normalfont\sffamily
\textbf{Aprendizagem contrastiva supervisionada}\par
Usar o rótulo fiscal para decidir quais descrições aproximar.
\par\vspace{1.5mm}
\renewcommand{\arraystretch}{1.35}
\begin{tabular}{@{}lp{.56\textwidth}l@{}}
\toprule
\textbf{Papel} & \textbf{Descrição ilustrativa} & \textbf{Classe}\\
\midrule
Âncora & Notebook Lenovo ThinkPad 512 GB & A\\
Positivo & Laptop ThinkPad Lenovo, SSD 512GB & A\\
Negativo & Monitor LCD para computador & B\\
\bottomrule
\end{tabular}\par
\par\vspace{1.5mm}
\alert{Aproximar a âncora do positivo e favorecer sua distinção dos negativos.}
\par\par\vspace{1.5mm}
{\small Exemplo didático: A/B não são enquadramentos fiscais reais.}
\fonte{Anggoro et al., seção 3.2 e Figura 2. Exemplo elaborado para esta apresentação.}
\end{frame}
\note{\small A âncora é uma descrição da base e o positivo é outra transação com o mesmo código. Os textos não precisam ser idênticos nem o produto precisa ser exatamente o mesmo. O método usa os rótulos para formar pares sem gerar paráfrases artificiais que possam eliminar informação técnica. O artigo experimenta um ou três positivos por âncora, formando pares adicionais. Os negativos serão obtidos a partir do lote pela MNR.\par [Sources] Anggoro et al., seção 3.2 e Figura 2. Exemplo elaborado para esta apresentação.}

\renewcommand{\secao}{CONCEITOS}
% Slide 7; 80 segundos
\begin{frame}[t]{MNR compara o positivo com os demais}
\normalfont\sffamily
\textbf{MNR: Multiple Negative Ranking Loss}\par
Função de perda já existente, aplicada ao ajuste do SBERT.
\par\vspace{1.5mm}
\begin{center}
\renewcommand{\arraystretch}{1.35}
\begin{tabular}{@{}lcc@{}}
\toprule
 & Positivo $b_1$ & Positivo $b_2$\\
\midrule
Âncora $a_1$ & \textcolor{Verde}{\textbf{par correto}} & negativo do lote\\
Âncora $a_2$ & negativo do lote & \textcolor{Verde}{\textbf{par correto}}\\
\bottomrule
\end{tabular}\par
\end{center}
\par\vspace{1.5mm}
\begin{itemize}
\item A perda favorece maior similaridade para o par correto.
\item Usa outros positivos do lote como negativos de cada âncora.
\item Similaridade por cosseno ou produto escalar.
\end{itemize}
\par\vspace{1mm}
{\small Cuidado: outro item do mesmo código pode virar um falso negativo.}
\fonte{Anggoro et al., seção 3.3 e Apêndice II. Falsos negativos: análise do apresentador.}
\end{frame}
\note{\small Para um lote com n pares, a âncora ai deve preferir bi aos demais bj. É uma competição de similaridades: não basta aproximar dois textos isoladamente. A MNR usa exemplos negativos do próprio lote, poupando seleção manual. A técnica não foi inventada por Anggoro: seu artigo referencia Henderson et al. (2017) e literatura posterior. A escolha dos pares pelo código é a aplicação supervisionada ao domínio. Apontar como análise crítica que a hipótese de outros pares serem negativos pode falhar se o lote contiver vários exemplos do mesmo código; não afirmar que o artigo documenta uma solução específica para isso.\par [Sources] Anggoro et al., seção 3.3 e Apêndice II; Henderson et al. (2017), referência do artigo. Ressalva de falsos negativos: análise metodológica do apresentador.}

\renewcommand{\secao}{METODOLOGIA}
% Slide 8; 75 segundos
\begin{frame}[t]{Há dois treinamentos separados}
\normalfont\sffamily
\begin{columns}[T,onlytextwidth]
\begin{column}{.47\textwidth}
\textcolor{Verde}{\textbf{1. Ajustar a representação}}\par\par\vspace{1.5mm}
\textbf{70\% das transações}\par\par\vspace{1.5mm}
Descrições + códigos\\[1mm]
$\downarrow$\\[1mm]
Pares positivos\\[1mm]
$\downarrow$\\[1mm]
SBERT treinado com MNR
\end{column}
\begin{column}{.47\textwidth}
\textcolor{Verde}{\textbf{2. Treinar o classificador}}\par\par\vspace{1.5mm}
\textbf{30\% das transações}\par\par\vspace{1.5mm}
SBERT ajustado gera vetores\\[1mm]
$\downarrow$\\[1mm]
SVM ou Random Forest\\[1mm]
$\downarrow$\\[1mm]
Validação estratificada: 5 folds
\end{column}
\end{columns}\par
\par\vspace{1.5mm}
\begin{block}{Uso em uma descrição nova}
Texto $\rightarrow$ SBERT ajustado $\rightarrow$ vetor $\rightarrow$ classificador $\rightarrow$ código
\end{block}
\fonte{Anggoro et al., seção 3.1, seção sobre classificadores e seção 5.1.}
\end{frame}
\note{\small O protocolo não é uma única divisão convencional de 70\% treino e 30\% teste de toda a cadeia. Setenta por cento ajustam o gerador de embeddings; os outros trinta por cento são destinados ao treino e à avaliação dos classificadores com cinco partições estratificadas. A estratificação preserva aproximadamente as proporções das classes, não torna a base balanceada. Na segunda etapa o codificador ajustado produz features; são os parâmetros do SVM ou da floresta que são treinados. SVM aprende fronteiras de decisão; Random Forest combina árvores.\par [Sources] Anggoro et al., seção 3.1, seção sobre classificadores e seção 5.1.}

\renewcommand{\secao}{METODOLOGIA}
% Slide 9; 55 segundos
\begin{frame}[t]{Os experimentos isolam escolhas do método}
\normalfont\sffamily
\renewcommand{\arraystretch}{1.35}
\begin{tabular}{@{}lp{.62\textwidth}@{}}
\toprule
\textbf{Escolha} & \textbf{Alternativas avaliadas}\\
\midrule
Codificador & BERT / DistilBERT (base, uncased)\\
Texto & Original / pré-processado\\
Pares positivos & 1 / 3 por âncora\\
Pooling & Média / máximo / [CLS]\\
Perda & MNR / MNR simétrica; cosseno / produto escalar\\
Treinamento & Lotes 8/32; épocas 10/30; taxas 0,00002 / 0,00005\\
\bottomrule
\end{tabular}\par
\par\vspace{1.5mm}
\textbf{Classificadores:} SVM e Random Forest, em configurações padrão.\par
\par\vspace{1.5mm}
{\small Execução com GPU; comparação com BERT/DistilBERT ajustados diretamente.}
\fonte{Anggoro et al., seções 4.2--4.3 e Tabelas III--VIII.}
\end{frame}
\note{\small Essa visão permite acompanhar quais componentes mudam. As tabelas seguintes são comparações de configurações, não um único experimento com todos os parâmetros possíveis combinados. MNR simétrica considera os dois sentidos do par. A escala foi avaliada em 20 e 50. O artigo não faz uma otimização adicional dos hiperparâmetros dos classificadores e usa recursos de cluster com GPU. A referência é transformer ajustado diretamente para prever classes, não BERT sem treinamento.\par [Sources] Anggoro et al., seções 4.2--4.3 e Tabelas III--VIII.}

\renewcommand{\secao}{AVALIAÇÃO}
% Slide 10; 40 segundos
\begin{frame}[t]{O F1 deve ser lido junto ao desbalanceamento}
\normalfont\sffamily
\begin{itemize}
\item \textbf{Precisão:} entre as previsões de uma classe, quantas estão corretas?
\item \textbf{Revocação:} entre os itens da classe, quantos foram recuperados?
\item \textbf{F1:} combina precisão e revocação.
\end{itemize}
\par\vspace{1.5mm}
\begin{columns}[T,onlytextwidth]
\begin{column}{.47\textwidth}
\textbf{F1 ponderado}\par\par\vspace{1.5mm}
Pesa cada classe pela sua frequência. É o foco das comparações do artigo.
\end{column}
\begin{column}{.47\textwidth}
\textbf{Macro-F1}\par\par\vspace{1.5mm}
Dá o mesmo peso a cada classe. Ajuda a avaliar classes menos frequentes.
\end{column}
\end{columns}\par
\par\vspace{1.5mm}
\alert{Um bom resultado agregado não garante bom desempenho em cada código.}
\fonte{Anggoro et al., seção 5.1 e Tabela III.}
\end{frame}
\note{\small O artigo relata kappa de Cohen e médias macro, micro e ponderada de precisão, revocação e F1. Aqui as comparações numéricas usarão F1 ponderado para manter consistência com o foco do artigo. Ele não deve ser interpretado como número de decisões fiscais auditadas corretamente. Em dados desbalanceados, resultados por classe e macro-F1 completam a leitura. Não reproduzir como equivalentes as médias de F1 e a acurácia de outro estudo.\par [Sources] Anggoro et al., seção 5.1 e Tabela III.}

\renewcommand{\secao}{RESULTADOS}
% Slide 11; 65 segundos
\begin{frame}[t]{O ajuste contrastivo supera a referência}
\normalfont\sffamily
\textbf{F1 ponderado}\par\par\vspace{1.5mm}
\renewcommand{\arraystretch}{1.4}
\begin{tabular}{@{}lcc@{}}
\toprule
\textbf{Método} & \textbf{Índia} & \textbf{EUA}\\
\midrule
BERT ajustado diretamente & 0,8406 & 0,7689\\
SBERT/BERT + SVM & 0,8620 & \textcolor{Verde}{\textbf{0,8013}}\\
SBERT/BERT + Random Forest & \textcolor{Verde}{\textbf{0,8656}} & 0,8005\\
\bottomrule
\end{tabular}\par
\par\vspace{1.5mm}
\begin{columns}[T,onlytextwidth]
\begin{column}{.47\textwidth}
{\LARGE\color{Verde}\textbf{+2,50 p.p.}}\par\par\vspace{1.5mm}
Melhor ganho na Índia.
\end{column}
\begin{column}{.47\textwidth}
{\LARGE\color{Verde}\textbf{+3,24 p.p.}}\par\par\vspace{1.5mm}
Melhor ganho nos EUA.
\end{column}
\end{columns}\par
\par\vspace{1.5mm}
{\small Ganhos calculados em relação ao BERT da Tabela III.}
\fonte{Anggoro et al., Tabela III; diferenças calculadas a partir dos valores publicados.}
\end{frame}
\note{\small Os ganhos são diferenças absolutas em pontos percentuais, não porcentagens relativas. Índia: 0,8656 menos 0,8406 = 0,0250. EUA: 0,8013 menos 0,7689 = 0,0324. A tabela mostra RF melhor na Índia e SVM ligeiramente melhor nos EUA. Não há vencedor universal entre os classificadores. A comparação sustenta o método nessas bases; não demonstra transferência entre países nem superioridade em toda tarefa de classificação fiscal. As tabelas de ablação contêm outras configurações e alguns resultados superiores, portanto evitar chamar esses valores de máximos globais do artigo.\par [Sources] Anggoro et al., Tabela III; diferenças calculadas a partir dos valores publicados.}

\renewcommand{\secao}{RESULTADOS}
% Slide 12; 60 segundos
\begin{frame}[t]{A limpeza remove informação útil}
\normalfont\sffamily
\textbf{SBERT/DistilBERT + SVM: F1 ponderado}\par\par\vspace{1.5mm}
\renewcommand{\arraystretch}{1.4}
\begin{tabular}{@{}lccc@{}}
\toprule
\textbf{Base} & \textbf{Original} & \textbf{Limpo} & \textbf{Variação}\\
\midrule
Índia & 0,8603 & 0,7790 & \textcolor{Terracota}{\textbf{$-$8,13 p.p.}}\\
EUA & 0,7953 & 0,7185 & \textcolor{Terracota}{\textbf{$-$7,68 p.p.}}\\
\bottomrule
\end{tabular}\par
\par\vspace{1.5mm}
\textbf{O que pode se perder?}\par\par\vspace{1.5mm}
Modelo, capacidade e especificações como \textbf{512 GB}.\par
\par\vspace{1.5mm}
\begin{itemize}
\item LIME: explica localmente quais termos influenciam uma previsão.
\item O exemplo do artigo indica relevância de atributos técnicos.
\end{itemize}
\fonte{Anggoro et al., seção 3.1, seção 5.3, Tabela IV e Figura 3.}
\end{frame}
\note{\small O pré-processamento inclui conversão para minúsculas, remoção de duplicatas, limpeza de informação alfanumérica e retenção de tokens com pelo menos dois caracteres. A redução de tokens foi 8\% na Índia e 19\% nos EUA. A Tabela IV mostra queda também com RF; os dois números exibidos são do SVM com DistilBERT, não da configuração BERT do slide anterior. LIME analisa uma previsão local por perturbações do texto. A Figura 3 ilustra que a informação 512 GB influencia a decisão. Isso não prova causalidade geral nem que toda limpeza de texto prejudica.\par [Sources] Anggoro et al., seção 3.1, seção 5.3, Tabela IV e Figura 3.}

\renewcommand{\secao}{RESULTADOS}
% Slide 13; 50 segundos
\begin{frame}[t]{Mais positivos enriquecem o ajuste}
\normalfont\sffamily
\textbf{SBERT/DistilBERT: F1 ponderado}\par\par\vspace{1.5mm}
\renewcommand{\arraystretch}{1.4}
\begin{tabular}{@{}llccc@{}}
\toprule
\textbf{Base} & \textbf{Modelo} & \textbf{1 positivo} & \textbf{3 positivos} & \textbf{Ganho}\\
\midrule
Índia & SVM & 0,8360 & 0,8603 & +2,43 p.p.\\
Índia & RF & 0,8410 & 0,8637 & +2,27 p.p.\\
EUA & SVM & 0,7735 & 0,7953 & +2,18 p.p.\\
EUA & RF & 0,7688 & 0,7958 & +2,70 p.p.\\
\bottomrule
\end{tabular}\par
\par\vspace{1.5mm}
\textbf{Interpretação:} mais descrições da mesma classe oferecem maior diversidade ao treinamento.\par
\par\vspace{1.5mm}
\alert{O benefício tem custo: mais pares e maior tempo de ajuste.}
\fonte{Anggoro et al., seção 3.2, seção 4.3.1 e Tabela V; ganhos calculados.}
\end{frame}
\note{\small Cada âncora recebe três positivos em vez de um, produzindo mais pares de treinamento. O ganho é consistente nos quatro cenários da Tabela V, entre 2,18 e 2,70 pontos percentuais. Não extrapolar para um número ilimitado de positivos; o artigo não demonstra esse comportamento. O mesmo código pode agrupar várias formas legítimas de descrever mercadorias, e essa diversidade é o sinal que o codificador aprende.\par [Sources] Anggoro et al., seção 3.2, seção 4.3.1 e Tabela V; ganhos calculados.}

\renewcommand{\secao}{RESULTADOS}
% Slide 14; 45 segundos
\begin{frame}[t]{O ganho também tem custo de armazenamento}
\normalfont\sffamily
\begin{itemize}
\item \textbf{Pooling:} resultados próximos; nenhuma opção domina todas as configurações.
\item \textbf{MNR simétrica e escala:} mudanças pequenas nos cenários avaliados.
\item \textbf{Treinamento:} 30 épocas e taxa 0,00002 favorecem os resultados; lotes maiores trazem ganhos modestos.
\end{itemize}
\par\vspace{1.5mm}
\textbf{Tamanho relatado na base indiana}\par\par\vspace{1.5mm}
\renewcommand{\arraystretch}{1.25}
\begin{tabular}{@{}lrr@{}}
\toprule
 & \textbf{F1 ponderado} & \textbf{Tamanho}\\
\midrule
BERT ajustado diretamente & 0,8406 & 418 MB\\
SBERT/DistilBERT + SVM & 0,8603 & 291 MB\\
SBERT/BERT + RF & 0,8656 & 631 MB\\
\bottomrule
\end{tabular}\par
\fonte{Anggoro et al., Tabelas III e VI--VIII.}
\end{frame}
\note{\small As Tabelas VI a VIII avaliam componentes. Não atribuir o ganho principal a uma única escolha de pooling. A tabela de tamanho usa a Tabela III: DistilBERT com SVM ocupa menos que BERT diretamente ajustado, mas mais que DistilBERT diretamente ajustado, que tem 255 MB. Modelo menor não implica necessariamente menor latência de toda a cadeia: também há o classificador. As comparações de épocas e lotes são experimentos próprios; não misturar números para inventar uma configuração ótima conjunta.\par [Sources] Anggoro et al., Tabelas III e VI--VIII.}

\renewcommand{\secao}{COMPLEMENTO}
% Slide 15; material complementar
\begin{frame}[t]{Como a MNR calcula a perda}
\normalfont\sffamily
\textbf{Em cada lote de pares positivos:}
\par\vspace{1.5mm}
\begin{enumerate}
\item Calcular a similaridade de cada âncora com todos os candidatos do lote.
\item Aplicar a escala e converter os escores em probabilidades (softmax).
\item Penalizar o modelo quando o positivo correto recebe baixa probabilidade.
\item Atualizar o codificador para reduzir essa perda.
\end{enumerate}
\par\vspace{1.5mm}
\textbf{MNR simétrica:} também avalia os pares no sentido inverso.
\par\par\vspace{1.5mm}
{\small A perda padrão é a entropia cruzada sobre os candidatos do lote.}
\fonte{Anggoro et al., seção 3.3, Tabela VII e Apêndice II; descrição didática do cálculo.}
\end{frame}
\note{\small A MNR padrão calcula entropia cruzada sobre candidatos do lote. As similaridades passam por uma escala e pela softmax. O rótulo esperado é a posição do positivo correspondente. A versão simétrica considera o sentido inverso. A escala avaliada no artigo foi 20 ou 50; não é uma regra fiscal.\par [Sources] Anggoro et al., seção 3.3, Tabela VII e Apêndice II; descrição didática do cálculo.}

\renewcommand{\secao}{REFERÊNCIAS}
% Slide 16; material complementar
\begin{frame}[t]{Fontes utilizadas}
\normalfont\sffamily
{\small
\textbf{Anggoro, A. W.; Corcoran, P.; De Widt, D.; Li, Y. (2025).}\par
\textit{Harmonized System Code Classification using Supervised Contrastive Learning with Sentence BERT and Multiple Negative Ranking Loss.}\par
Data Technologies and Applications, 59(2), 279--301.\par
\href{https://doi.org/10.1108/DTA-01-2024-0052}{DOI: 10.1108/DTA-01-2024-0052}. Versão dos autores, ORCA.
\par\vspace{1.5mm}
\textbf{Paula, W.; Jacob Junior, A. F. L.; Cortes, O. A. C.; Lobato, F. M.}\par
\textit{Avaliação de Large Language Models Para Classificação de Produtos da Nota Fiscal Eletrônica em Grupos da Nomenclatura Comum do Mercosul.}\par
Manuscrito local: artigo\_welyab\_pecs.pdf.
\par\vspace{1.5mm}
\textbf{Referências técnicas citadas por Anggoro}\par
Reimers e Gurevych (2019): SBERT. Henderson et al. (2017): aprendizagem com negativos. Ribeiro et al. (2016): LIME.
}
\fonte{PDFs locais de Anggoro et al. e Welyab et al.; bibliografia de Anggoro.}
\end{frame}
\note{\small O artigo de Anggoro é a fonte primária dos resultados e conceitos apresentados; o PDF de Welyab fundamenta o contexto e a comparação metodológica. As referências técnicas são apresentadas conforme a bibliografia de Anggoro, sem sugerir experimentos próprios.\par [Sources] PDFs locais de Anggoro et al. e Welyab et al.; bibliografia de Anggoro.}
\end{document}
